% F083｜首页以实际指数分式曲线呈现原点对称。
\begin{statementbox}
% Formula-based samples; see F083 checks/build_tikz_diagrams.py.
\begin{center}\begin{tikzpicture}[font=\small,>=stealth,line join=round]
\draw[gray!40,rounded corners=3pt] (0,0) rectangle (14.2,8.4);
\node at (7.1,7.75) {奇2：$y=\dfrac{a^x-1}{a^x+1}$};
\begin{scope}[shift={(4.400000,3.900000)},x=0.850000cm,y=1.730769cm]
\draw[gray,dotted] (-4,-1)--(4,-1);
\node[anchor=east] at (-4-.08/0.850000,-1) {$-1$};
\draw[gray,dotted] (-4,1)--(4,1);
\node[anchor=east] at (-4-.08/0.850000,1) {$1$};
\draw[->,gray] (-4,0)--(4,0) node[right] {$x$};
\draw[->,gray] (0,-1.3)--(0,1.3) node[above] {$y$};
\draw[gray] (-1,-.035/1.730769)--(-1,.035/1.730769);
\node[anchor=north] at (-1,-.07/1.730769) {$-1$};
\draw[gray] (1,-.035/1.730769)--(1,.035/1.730769);
\node[anchor=north] at (1,-.07/1.730769) {$1$};
\begin{scope}
\clip (-4,-1.3) rectangle (4,1.3);
\draw[blue,thick,solid] plot coordinates {(-4.0000000,-0.8823529) (-3.9866667,-0.8813254) (-3.9733333,-0.8802895) (-3.9600000,-0.8792451) (-3.9466667,-0.8781922) (-3.9333333,-0.8771307) (-3.9200000,-0.8760606) (-3.9066667,-0.8749818) (-3.8933333,-0.8738942) (-3.8800000,-0.8727978) (-3.8666667,-0.8716926) (-3.8533333,-0.8705783) (-3.8400000,-0.8694551) (-3.8266667,-0.8683228) (-3.8133333,-0.8671814) (-3.8000000,-0.8660308) (-3.7866667,-0.8648710) (-3.7733333,-0.8637019) (-3.7600000,-0.8625233) (-3.7466667,-0.8613354) (-3.7333333,-0.8601380) (-3.7200000,-0.8589310) (-3.7066667,-0.8577144) (-3.6933333,-0.8564881) (-3.6800000,-0.8552520) (-3.6666667,-0.8540062) (-3.6533333,-0.8527504) (-3.6400000,-0.8514848) (-3.6266667,-0.8502091) (-3.6133333,-0.8489234) (-3.6000000,-0.8476276) (-3.5866667,-0.8463215) (-3.5733333,-0.8450052) (-3.5600000,-0.8436786) (-3.5466667,-0.8423416) (-3.5333333,-0.8409942) (-3.5200000,-0.8396362) (-3.5066667,-0.8382677) (-3.4933333,-0.8368885) (-3.4800000,-0.8354986) (-3.4666667,-0.8340979) (-3.4533333,-0.8326864) (-3.4400000,-0.8312640) (-3.4266667,-0.8298306) (-3.4133333,-0.8283862) (-3.4000000,-0.8269307) (-3.3866667,-0.8254640) (-3.3733333,-0.8239861) (-3.3600000,-0.8224969) (-3.3466667,-0.8209963) (-3.3333333,-0.8194843) (-3.3200000,-0.8179608) (-3.3066667,-0.8164257) (-3.2933333,-0.8148791) (-3.2800000,-0.8133207) (-3.2666667,-0.8117505) (-3.2533333,-0.8101686) (-3.2400000,-0.8085747) (-3.2266667,-0.8069689) (-3.2133333,-0.8053511) (-3.2000000,-0.8037212) (-3.1866667,-0.8020791) (-3.1733333,-0.8004248) (-3.1600000,-0.7987583) (-3.1466667,-0.7970794) (-3.1333333,-0.7953880) (-3.1200000,-0.7936842) (-3.1066667,-0.7919679) (-3.0933333,-0.7902389) (-3.0800000,-0.7884973) (-3.0666667,-0.7867429) (-3.0533333,-0.7849758) (-3.0400000,-0.7831957) (-3.0266667,-0.7814028) (-3.0133333,-0.7795968) (-3.0000000,-0.7777778) (-2.9866667,-0.7759456) (-2.9733333,-0.7741003) (-2.9600000,-0.7722417) (-2.9466667,-0.7703698) (-2.9333333,-0.7684846) (-2.9200000,-0.7665859) (-2.9066667,-0.7646737) (-2.8933333,-0.7627479) (-2.8800000,-0.7608085) (-2.8666667,-0.7588554) (-2.8533333,-0.7568886) (-2.8400000,-0.7549080) (-2.8266667,-0.7529135) (-2.8133333,-0.7509051) (-2.8000000,-0.7488827) (-2.7866667,-0.7468462) (-2.7733333,-0.7447957) (-2.7600000,-0.7427309) (-2.7466667,-0.7406520) (-2.7333333,-0.7385588) (-2.7200000,-0.7364512) (-2.7066667,-0.7343293) (-2.6933333,-0.7321929) (-2.6800000,-0.7300420) (-2.6666667,-0.7278765) (-2.6533333,-0.7256964) (-2.6400000,-0.7235017) (-2.6266667,-0.7212922) (-2.6133333,-0.7190679) (-2.6000000,-0.7168288) (-2.5866667,-0.7145749) (-2.5733333,-0.7123060) (-2.5600000,-0.7100221) (-2.5466667,-0.7077232) (-2.5333333,-0.7054091) (-2.5200000,-0.7030800) (-2.5066667,-0.7007357) (-2.4933333,-0.6983761) (-2.4800000,-0.6960013) (-2.4666667,-0.6936111) (-2.4533333,-0.6912056) (-2.4400000,-0.6887846) (-2.4266667,-0.6863482) (-2.4133333,-0.6838963) (-2.4000000,-0.6814288) (-2.3866667,-0.6789458) (-2.3733333,-0.6764471) (-2.3600000,-0.6739327) (-2.3466667,-0.6714027) (-2.3333333,-0.6688569) (-2.3200000,-0.6662953) (-2.3066667,-0.6637178) (-2.2933333,-0.6611246) (-2.2800000,-0.6585154) (-2.2666667,-0.6558903) (-2.2533333,-0.6532492) (-2.2400000,-0.6505922) (-2.2266667,-0.6479191) (-2.2133333,-0.6452300) (-2.2000000,-0.6425248) (-2.1866667,-0.6398035) (-2.1733333,-0.6370660) (-2.1600000,-0.6343124) (-2.1466667,-0.6315425) (-2.1333333,-0.6287565) (-2.1200000,-0.6259542) (-2.1066667,-0.6231357) (-2.0933333,-0.6203009) (-2.0800000,-0.6174498) (-2.0666667,-0.6145824) (-2.0533333,-0.6116986) (-2.0400000,-0.6087985) (-2.0266667,-0.6058821) (-2.0133333,-0.6029492) (-2.0000000,-0.6000000) (-1.9866667,-0.5970344) (-1.9733333,-0.5940523) (-1.9600000,-0.5910539) (-1.9466667,-0.5880390) (-1.9333333,-0.5850077) (-1.9200000,-0.5819599) (-1.9066667,-0.5788958) (-1.8933333,-0.5758152) (-1.8800000,-0.5727181) (-1.8666667,-0.5696046) (-1.8533333,-0.5664747) (-1.8400000,-0.5633283) (-1.8266667,-0.5601656) (-1.8133333,-0.5569864) (-1.8000000,-0.5537908) (-1.7866667,-0.5505788) (-1.7733333,-0.5473504) (-1.7600000,-0.5441056) (-1.7466667,-0.5408445) (-1.7333333,-0.5375671) (-1.7200000,-0.5342733) (-1.7066667,-0.5309632) (-1.6933333,-0.5276369) (-1.6800000,-0.5242942) (-1.6666667,-0.5209354) (-1.6533333,-0.5175603) (-1.6400000,-0.5141691) (-1.6266667,-0.5107617) (-1.6133333,-0.5073381) (-1.6000000,-0.5038985) (-1.5866667,-0.5004428) (-1.5733333,-0.4969711) (-1.5600000,-0.4934835) (-1.5466667,-0.4899798) (-1.5333333,-0.4864603) (-1.5200000,-0.4829250) (-1.5066667,-0.4793738) (-1.4933333,-0.4758068) (-1.4800000,-0.4722241) (-1.4666667,-0.4686258) (-1.4533333,-0.4650118) (-1.4400000,-0.4613823) (-1.4266667,-0.4577372) (-1.4133333,-0.4540767) (-1.4000000,-0.4504009) (-1.3866667,-0.4467096) (-1.3733333,-0.4430031) (-1.3600000,-0.4392814) (-1.3466667,-0.4355446) (-1.3333333,-0.4317927) (-1.3200000,-0.4280258) (-1.3066667,-0.4242439) (-1.2933333,-0.4204472) (-1.2800000,-0.4166358) (-1.2666667,-0.4128096) (-1.2533333,-0.4089688) (-1.2400000,-0.4051134) (-1.2266667,-0.4012436) (-1.2133333,-0.3973594) (-1.2000000,-0.3934609) (-1.1866667,-0.3895482) (-1.1733333,-0.3856214) (-1.1600000,-0.3816806) (-1.1466667,-0.3777259) (-1.1333333,-0.3737573) (-1.1200000,-0.3697750) (-1.1066667,-0.3657790) (-1.0933333,-0.3617696) (-1.0800000,-0.3577467) (-1.0666667,-0.3537105) (-1.0533333,-0.3496610) (-1.0400000,-0.3455985) (-1.0266667,-0.3415229) (-1.0133333,-0.3374345) (-1.0000000,-0.3333333) (-0.9866667,-0.3292195) (-0.9733333,-0.3250931) (-0.9600000,-0.3209543) (-0.9466667,-0.3168032) (-0.9333333,-0.3126399) (-0.9200000,-0.3084646) (-0.9066667,-0.3042774) (-0.8933333,-0.3000784) (-0.8800000,-0.2958677) (-0.8666667,-0.2916455) (-0.8533333,-0.2874119) (-0.8400000,-0.2831670) (-0.8266667,-0.2789110) (-0.8133333,-0.2746440) (-0.8000000,-0.2703662) (-0.7866667,-0.2660777) (-0.7733333,-0.2617786) (-0.7600000,-0.2574691) (-0.7466667,-0.2531493) (-0.7333333,-0.2488195) (-0.7200000,-0.2444796) (-0.7066667,-0.2401299) (-0.6933333,-0.2357706) (-0.6800000,-0.2314018) (-0.6666667,-0.2270236) (-0.6533333,-0.2226362) (-0.6400000,-0.2182398) (-0.6266667,-0.2138345) (-0.6133333,-0.2094204) (-0.6000000,-0.2049979) (-0.5866667,-0.2005669) (-0.5733333,-0.1961278) (-0.5600000,-0.1916805) (-0.5466667,-0.1872254) (-0.5333333,-0.1827626) (-0.5200000,-0.1782922) (-0.5066667,-0.1738145) (-0.4933333,-0.1693295) (-0.4800000,-0.1648375) (-0.4666667,-0.1603387) (-0.4533333,-0.1558332) (-0.4400000,-0.1513213) (-0.4266667,-0.1468030) (-0.4133333,-0.1422785) (-0.4000000,-0.1377481) (-0.3866667,-0.1332120) (-0.3733333,-0.1286702) (-0.3600000,-0.1241231) (-0.3466667,-0.1195707) (-0.3333333,-0.1150133) (-0.3200000,-0.1104511) (-0.3066667,-0.1058842) (-0.2933333,-0.1013128) (-0.2800000,-0.0967371) (-0.2666667,-0.0921574) (-0.2533333,-0.0875737) (-0.2400000,-0.0829864) (-0.2266667,-0.0783955) (-0.2133333,-0.0738013) (-0.2000000,-0.0692039) (-0.1866667,-0.0646036) (-0.1733333,-0.0600006) (-0.1600000,-0.0553950) (-0.1466667,-0.0507871) (-0.1333333,-0.0461769) (-0.1200000,-0.0415649) (-0.1066667,-0.0369510) (-0.0933333,-0.0323356) (-0.0800000,-0.0277188) (-0.0666667,-0.0231008) (-0.0533333,-0.0184818) (-0.0400000,-0.0138621) (-0.0266667,-0.0092417) (-0.0133333,-0.0046209) (0.0000000,0.0000000) (0.0133333,0.0046209) (0.0266667,0.0092417) (0.0400000,0.0138621) (0.0533333,0.0184818) (0.0666667,0.0231008) (0.0800000,0.0277188) (0.0933333,0.0323356) (0.1066667,0.0369510) (0.1200000,0.0415649) (0.1333333,0.0461769) (0.1466667,0.0507871) (0.1600000,0.0553950) (0.1733333,0.0600006) (0.1866667,0.0646036) (0.2000000,0.0692039) (0.2133333,0.0738013) (0.2266667,0.0783955) (0.2400000,0.0829864) (0.2533333,0.0875737) (0.2666667,0.0921574) (0.2800000,0.0967371) (0.2933333,0.1013128) (0.3066667,0.1058842) (0.3200000,0.1104511) (0.3333333,0.1150133) (0.3466667,0.1195707) (0.3600000,0.1241231) (0.3733333,0.1286702) (0.3866667,0.1332120) (0.4000000,0.1377481) (0.4133333,0.1422785) (0.4266667,0.1468030) (0.4400000,0.1513213) (0.4533333,0.1558332) (0.4666667,0.1603387) (0.4800000,0.1648375) (0.4933333,0.1693295) (0.5066667,0.1738145) (0.5200000,0.1782922) (0.5333333,0.1827626) (0.5466667,0.1872254) (0.5600000,0.1916805) (0.5733333,0.1961278) (0.5866667,0.2005669) (0.6000000,0.2049979) (0.6133333,0.2094204) (0.6266667,0.2138345) (0.6400000,0.2182398) (0.6533333,0.2226362) (0.6666667,0.2270236) (0.6800000,0.2314018) (0.6933333,0.2357706) (0.7066667,0.2401299) (0.7200000,0.2444796) (0.7333333,0.2488195) (0.7466667,0.2531493) (0.7600000,0.2574691) (0.7733333,0.2617786) (0.7866667,0.2660777) (0.8000000,0.2703662) (0.8133333,0.2746440) (0.8266667,0.2789110) (0.8400000,0.2831670) (0.8533333,0.2874119) (0.8666667,0.2916455) (0.8800000,0.2958677) (0.8933333,0.3000784) (0.9066667,0.3042774) (0.9200000,0.3084646) (0.9333333,0.3126399) (0.9466667,0.3168032) (0.9600000,0.3209543) (0.9733333,0.3250931) (0.9866667,0.3292195) (1.0000000,0.3333333) (1.0133333,0.3374345) (1.0266667,0.3415229) (1.0400000,0.3455985) (1.0533333,0.3496610) (1.0666667,0.3537105) (1.0800000,0.3577467) (1.0933333,0.3617696) (1.1066667,0.3657790) (1.1200000,0.3697750) (1.1333333,0.3737573) (1.1466667,0.3777259) (1.1600000,0.3816806) (1.1733333,0.3856214) (1.1866667,0.3895482) (1.2000000,0.3934609) (1.2133333,0.3973594) (1.2266667,0.4012436) (1.2400000,0.4051134) (1.2533333,0.4089688) (1.2666667,0.4128096) (1.2800000,0.4166358) (1.2933333,0.4204472) (1.3066667,0.4242439) (1.3200000,0.4280258) (1.3333333,0.4317927) (1.3466667,0.4355446) (1.3600000,0.4392814) (1.3733333,0.4430031) (1.3866667,0.4467096) (1.4000000,0.4504009) (1.4133333,0.4540767) (1.4266667,0.4577372) (1.4400000,0.4613823) (1.4533333,0.4650118) (1.4666667,0.4686258) (1.4800000,0.4722241) (1.4933333,0.4758068) (1.5066667,0.4793738) (1.5200000,0.4829250) (1.5333333,0.4864603) (1.5466667,0.4899798) (1.5600000,0.4934835) (1.5733333,0.4969711) (1.5866667,0.5004428) (1.6000000,0.5038985) (1.6133333,0.5073381) (1.6266667,0.5107617) (1.6400000,0.5141691) (1.6533333,0.5175603) (1.6666667,0.5209354) (1.6800000,0.5242942) (1.6933333,0.5276369) (1.7066667,0.5309632) (1.7200000,0.5342733) (1.7333333,0.5375671) (1.7466667,0.5408445) (1.7600000,0.5441056) (1.7733333,0.5473504) (1.7866667,0.5505788) (1.8000000,0.5537908) (1.8133333,0.5569864) (1.8266667,0.5601656) (1.8400000,0.5633283) (1.8533333,0.5664747) (1.8666667,0.5696046) (1.8800000,0.5727181) (1.8933333,0.5758152) (1.9066667,0.5788958) (1.9200000,0.5819599) (1.9333333,0.5850077) (1.9466667,0.5880390) (1.9600000,0.5910539) (1.9733333,0.5940523) (1.9866667,0.5970344) (2.0000000,0.6000000) (2.0133333,0.6029492) (2.0266667,0.6058821) (2.0400000,0.6087985) (2.0533333,0.6116986) (2.0666667,0.6145824) (2.0800000,0.6174498) (2.0933333,0.6203009) (2.1066667,0.6231357) (2.1200000,0.6259542) (2.1333333,0.6287565) (2.1466667,0.6315425) (2.1600000,0.6343124) (2.1733333,0.6370660) (2.1866667,0.6398035) (2.2000000,0.6425248) (2.2133333,0.6452300) (2.2266667,0.6479191) (2.2400000,0.6505922) (2.2533333,0.6532492) (2.2666667,0.6558903) (2.2800000,0.6585154) (2.2933333,0.6611246) (2.3066667,0.6637178) (2.3200000,0.6662953) (2.3333333,0.6688569) (2.3466667,0.6714027) (2.3600000,0.6739327) (2.3733333,0.6764471) (2.3866667,0.6789458) (2.4000000,0.6814288) (2.4133333,0.6838963) (2.4266667,0.6863482) (2.4400000,0.6887846) (2.4533333,0.6912056) (2.4666667,0.6936111) (2.4800000,0.6960013) (2.4933333,0.6983761) (2.5066667,0.7007357) (2.5200000,0.7030800) (2.5333333,0.7054091) (2.5466667,0.7077232) (2.5600000,0.7100221) (2.5733333,0.7123060) (2.5866667,0.7145749) (2.6000000,0.7168288) (2.6133333,0.7190679) (2.6266667,0.7212922) (2.6400000,0.7235017) (2.6533333,0.7256964) (2.6666667,0.7278765) (2.6800000,0.7300420) (2.6933333,0.7321929) (2.7066667,0.7343293) (2.7200000,0.7364512) (2.7333333,0.7385588) (2.7466667,0.7406520) (2.7600000,0.7427309) (2.7733333,0.7447957) (2.7866667,0.7468462) (2.8000000,0.7488827) (2.8133333,0.7509051) (2.8266667,0.7529135) (2.8400000,0.7549080) (2.8533333,0.7568886) (2.8666667,0.7588554) (2.8800000,0.7608085) (2.8933333,0.7627479) (2.9066667,0.7646737) (2.9200000,0.7665859) (2.9333333,0.7684846) (2.9466667,0.7703698) (2.9600000,0.7722417) (2.9733333,0.7741003) (2.9866667,0.7759456) (3.0000000,0.7777778) (3.0133333,0.7795968) (3.0266667,0.7814028) (3.0400000,0.7831957) (3.0533333,0.7849758) (3.0666667,0.7867429) (3.0800000,0.7884973) (3.0933333,0.7902389) (3.1066667,0.7919679) (3.1200000,0.7936842) (3.1333333,0.7953880) (3.1466667,0.7970794) (3.1600000,0.7987583) (3.1733333,0.8004248) (3.1866667,0.8020791) (3.2000000,0.8037212) (3.2133333,0.8053511) (3.2266667,0.8069689) (3.2400000,0.8085747) (3.2533333,0.8101686) (3.2666667,0.8117505) (3.2800000,0.8133207) (3.2933333,0.8148791) (3.3066667,0.8164257) (3.3200000,0.8179608) (3.3333333,0.8194843) (3.3466667,0.8209963) (3.3600000,0.8224969) (3.3733333,0.8239861) (3.3866667,0.8254640) (3.4000000,0.8269307) (3.4133333,0.8283862) (3.4266667,0.8298306) (3.4400000,0.8312640) (3.4533333,0.8326864) (3.4666667,0.8340979) (3.4800000,0.8354986) (3.4933333,0.8368885) (3.5066667,0.8382677) (3.5200000,0.8396362) (3.5333333,0.8409942) (3.5466667,0.8423416) (3.5600000,0.8436786) (3.5733333,0.8450052) (3.5866667,0.8463215) (3.6000000,0.8476276) (3.6133333,0.8489234) (3.6266667,0.8502091) (3.6400000,0.8514848) (3.6533333,0.8527504) (3.6666667,0.8540062) (3.6800000,0.8552520) (3.6933333,0.8564881) (3.7066667,0.8577144) (3.7200000,0.8589310) (3.7333333,0.8601380) (3.7466667,0.8613354) (3.7600000,0.8625233) (3.7733333,0.8637019) (3.7866667,0.8648710) (3.8000000,0.8660308) (3.8133333,0.8671814) (3.8266667,0.8683228) (3.8400000,0.8694551) (3.8533333,0.8705783) (3.8666667,0.8716926) (3.8800000,0.8727978) (3.8933333,0.8738942) (3.9066667,0.8749818) (3.9200000,0.8760606) (3.9333333,0.8771307) (3.9466667,0.8781922) (3.9600000,0.8792451) (3.9733333,0.8802895) (3.9866667,0.8813254) (4.0000000,0.8823529)};
\draw[orange!85!black,thick,dashed] plot coordinates {(-4.0000000,0.8823529) (-3.9866667,0.8813254) (-3.9733333,0.8802895) (-3.9600000,0.8792451) (-3.9466667,0.8781922) (-3.9333333,0.8771307) (-3.9200000,0.8760606) (-3.9066667,0.8749818) (-3.8933333,0.8738942) (-3.8800000,0.8727978) (-3.8666667,0.8716926) (-3.8533333,0.8705783) (-3.8400000,0.8694551) (-3.8266667,0.8683228) (-3.8133333,0.8671814) (-3.8000000,0.8660308) (-3.7866667,0.8648710) (-3.7733333,0.8637019) (-3.7600000,0.8625233) (-3.7466667,0.8613354) (-3.7333333,0.8601380) (-3.7200000,0.8589310) (-3.7066667,0.8577144) (-3.6933333,0.8564881) (-3.6800000,0.8552520) (-3.6666667,0.8540062) (-3.6533333,0.8527504) (-3.6400000,0.8514848) (-3.6266667,0.8502091) (-3.6133333,0.8489234) (-3.6000000,0.8476276) (-3.5866667,0.8463215) (-3.5733333,0.8450052) (-3.5600000,0.8436786) (-3.5466667,0.8423416) (-3.5333333,0.8409942) (-3.5200000,0.8396362) (-3.5066667,0.8382677) (-3.4933333,0.8368885) (-3.4800000,0.8354986) (-3.4666667,0.8340979) (-3.4533333,0.8326864) (-3.4400000,0.8312640) (-3.4266667,0.8298306) (-3.4133333,0.8283862) (-3.4000000,0.8269307) (-3.3866667,0.8254640) (-3.3733333,0.8239861) (-3.3600000,0.8224969) (-3.3466667,0.8209963) (-3.3333333,0.8194843) (-3.3200000,0.8179608) (-3.3066667,0.8164257) (-3.2933333,0.8148791) (-3.2800000,0.8133207) (-3.2666667,0.8117505) (-3.2533333,0.8101686) (-3.2400000,0.8085747) (-3.2266667,0.8069689) (-3.2133333,0.8053511) (-3.2000000,0.8037212) (-3.1866667,0.8020791) (-3.1733333,0.8004248) (-3.1600000,0.7987583) (-3.1466667,0.7970794) (-3.1333333,0.7953880) (-3.1200000,0.7936842) (-3.1066667,0.7919679) (-3.0933333,0.7902389) (-3.0800000,0.7884973) (-3.0666667,0.7867429) (-3.0533333,0.7849758) (-3.0400000,0.7831957) (-3.0266667,0.7814028) (-3.0133333,0.7795968) (-3.0000000,0.7777778) (-2.9866667,0.7759456) (-2.9733333,0.7741003) (-2.9600000,0.7722417) (-2.9466667,0.7703698) (-2.9333333,0.7684846) (-2.9200000,0.7665859) (-2.9066667,0.7646737) (-2.8933333,0.7627479) (-2.8800000,0.7608085) (-2.8666667,0.7588554) (-2.8533333,0.7568886) (-2.8400000,0.7549080) (-2.8266667,0.7529135) (-2.8133333,0.7509051) (-2.8000000,0.7488827) (-2.7866667,0.7468462) (-2.7733333,0.7447957) (-2.7600000,0.7427309) (-2.7466667,0.7406520) (-2.7333333,0.7385588) (-2.7200000,0.7364512) (-2.7066667,0.7343293) (-2.6933333,0.7321929) (-2.6800000,0.7300420) (-2.6666667,0.7278765) (-2.6533333,0.7256964) (-2.6400000,0.7235017) (-2.6266667,0.7212922) (-2.6133333,0.7190679) (-2.6000000,0.7168288) (-2.5866667,0.7145749) (-2.5733333,0.7123060) (-2.5600000,0.7100221) (-2.5466667,0.7077232) (-2.5333333,0.7054091) (-2.5200000,0.7030800) (-2.5066667,0.7007357) (-2.4933333,0.6983761) (-2.4800000,0.6960013) (-2.4666667,0.6936111) (-2.4533333,0.6912056) (-2.4400000,0.6887846) (-2.4266667,0.6863482) (-2.4133333,0.6838963) (-2.4000000,0.6814288) (-2.3866667,0.6789458) (-2.3733333,0.6764471) (-2.3600000,0.6739327) (-2.3466667,0.6714027) (-2.3333333,0.6688569) (-2.3200000,0.6662953) (-2.3066667,0.6637178) (-2.2933333,0.6611246) (-2.2800000,0.6585154) (-2.2666667,0.6558903) (-2.2533333,0.6532492) (-2.2400000,0.6505922) (-2.2266667,0.6479191) (-2.2133333,0.6452300) (-2.2000000,0.6425248) (-2.1866667,0.6398035) (-2.1733333,0.6370660) (-2.1600000,0.6343124) (-2.1466667,0.6315425) (-2.1333333,0.6287565) (-2.1200000,0.6259542) (-2.1066667,0.6231357) (-2.0933333,0.6203009) (-2.0800000,0.6174498) (-2.0666667,0.6145824) (-2.0533333,0.6116986) (-2.0400000,0.6087985) (-2.0266667,0.6058821) (-2.0133333,0.6029492) (-2.0000000,0.6000000) (-1.9866667,0.5970344) (-1.9733333,0.5940523) (-1.9600000,0.5910539) (-1.9466667,0.5880390) (-1.9333333,0.5850077) (-1.9200000,0.5819599) (-1.9066667,0.5788958) (-1.8933333,0.5758152) (-1.8800000,0.5727181) (-1.8666667,0.5696046) (-1.8533333,0.5664747) (-1.8400000,0.5633283) (-1.8266667,0.5601656) (-1.8133333,0.5569864) (-1.8000000,0.5537908) (-1.7866667,0.5505788) (-1.7733333,0.5473504) (-1.7600000,0.5441056) (-1.7466667,0.5408445) (-1.7333333,0.5375671) (-1.7200000,0.5342733) (-1.7066667,0.5309632) (-1.6933333,0.5276369) (-1.6800000,0.5242942) (-1.6666667,0.5209354) (-1.6533333,0.5175603) (-1.6400000,0.5141691) (-1.6266667,0.5107617) (-1.6133333,0.5073381) (-1.6000000,0.5038985) (-1.5866667,0.5004428) (-1.5733333,0.4969711) (-1.5600000,0.4934835) (-1.5466667,0.4899798) (-1.5333333,0.4864603) (-1.5200000,0.4829250) (-1.5066667,0.4793738) (-1.4933333,0.4758068) (-1.4800000,0.4722241) (-1.4666667,0.4686258) (-1.4533333,0.4650118) (-1.4400000,0.4613823) (-1.4266667,0.4577372) (-1.4133333,0.4540767) (-1.4000000,0.4504009) (-1.3866667,0.4467096) (-1.3733333,0.4430031) (-1.3600000,0.4392814) (-1.3466667,0.4355446) (-1.3333333,0.4317927) (-1.3200000,0.4280258) (-1.3066667,0.4242439) (-1.2933333,0.4204472) (-1.2800000,0.4166358) (-1.2666667,0.4128096) (-1.2533333,0.4089688) (-1.2400000,0.4051134) (-1.2266667,0.4012436) (-1.2133333,0.3973594) (-1.2000000,0.3934609) (-1.1866667,0.3895482) (-1.1733333,0.3856214) (-1.1600000,0.3816806) (-1.1466667,0.3777259) (-1.1333333,0.3737573) (-1.1200000,0.3697750) (-1.1066667,0.3657790) (-1.0933333,0.3617696) (-1.0800000,0.3577467) (-1.0666667,0.3537105) (-1.0533333,0.3496610) (-1.0400000,0.3455985) (-1.0266667,0.3415229) (-1.0133333,0.3374345) (-1.0000000,0.3333333) (-0.9866667,0.3292195) (-0.9733333,0.3250931) (-0.9600000,0.3209543) (-0.9466667,0.3168032) (-0.9333333,0.3126399) (-0.9200000,0.3084646) (-0.9066667,0.3042774) (-0.8933333,0.3000784) (-0.8800000,0.2958677) (-0.8666667,0.2916455) (-0.8533333,0.2874119) (-0.8400000,0.2831670) (-0.8266667,0.2789110) (-0.8133333,0.2746440) (-0.8000000,0.2703662) (-0.7866667,0.2660777) (-0.7733333,0.2617786) (-0.7600000,0.2574691) (-0.7466667,0.2531493) (-0.7333333,0.2488195) (-0.7200000,0.2444796) (-0.7066667,0.2401299) (-0.6933333,0.2357706) (-0.6800000,0.2314018) (-0.6666667,0.2270236) (-0.6533333,0.2226362) (-0.6400000,0.2182398) (-0.6266667,0.2138345) (-0.6133333,0.2094204) (-0.6000000,0.2049979) (-0.5866667,0.2005669) (-0.5733333,0.1961278) (-0.5600000,0.1916805) (-0.5466667,0.1872254) (-0.5333333,0.1827626) (-0.5200000,0.1782922) (-0.5066667,0.1738145) (-0.4933333,0.1693295) (-0.4800000,0.1648375) (-0.4666667,0.1603387) (-0.4533333,0.1558332) (-0.4400000,0.1513213) (-0.4266667,0.1468030) (-0.4133333,0.1422785) (-0.4000000,0.1377481) (-0.3866667,0.1332120) (-0.3733333,0.1286702) (-0.3600000,0.1241231) (-0.3466667,0.1195707) (-0.3333333,0.1150133) (-0.3200000,0.1104511) (-0.3066667,0.1058842) (-0.2933333,0.1013128) (-0.2800000,0.0967371) (-0.2666667,0.0921574) (-0.2533333,0.0875737) (-0.2400000,0.0829864) (-0.2266667,0.0783955) (-0.2133333,0.0738013) (-0.2000000,0.0692039) (-0.1866667,0.0646036) (-0.1733333,0.0600006) (-0.1600000,0.0553950) (-0.1466667,0.0507871) (-0.1333333,0.0461769) (-0.1200000,0.0415649) (-0.1066667,0.0369510) (-0.0933333,0.0323356) (-0.0800000,0.0277188) (-0.0666667,0.0231008) (-0.0533333,0.0184818) (-0.0400000,0.0138621) (-0.0266667,0.0092417) (-0.0133333,0.0046209) (0.0000000,-0.0000000) (0.0133333,-0.0046209) (0.0266667,-0.0092417) (0.0400000,-0.0138621) (0.0533333,-0.0184818) (0.0666667,-0.0231008) (0.0800000,-0.0277188) (0.0933333,-0.0323356) (0.1066667,-0.0369510) (0.1200000,-0.0415649) (0.1333333,-0.0461769) (0.1466667,-0.0507871) (0.1600000,-0.0553950) (0.1733333,-0.0600006) (0.1866667,-0.0646036) (0.2000000,-0.0692039) (0.2133333,-0.0738013) (0.2266667,-0.0783955) (0.2400000,-0.0829864) (0.2533333,-0.0875737) (0.2666667,-0.0921574) (0.2800000,-0.0967371) (0.2933333,-0.1013128) (0.3066667,-0.1058842) (0.3200000,-0.1104511) (0.3333333,-0.1150133) (0.3466667,-0.1195707) (0.3600000,-0.1241231) (0.3733333,-0.1286702) (0.3866667,-0.1332120) (0.4000000,-0.1377481) (0.4133333,-0.1422785) (0.4266667,-0.1468030) (0.4400000,-0.1513213) (0.4533333,-0.1558332) (0.4666667,-0.1603387) (0.4800000,-0.1648375) (0.4933333,-0.1693295) (0.5066667,-0.1738145) (0.5200000,-0.1782922) (0.5333333,-0.1827626) (0.5466667,-0.1872254) (0.5600000,-0.1916805) (0.5733333,-0.1961278) (0.5866667,-0.2005669) (0.6000000,-0.2049979) (0.6133333,-0.2094204) (0.6266667,-0.2138345) (0.6400000,-0.2182398) (0.6533333,-0.2226362) (0.6666667,-0.2270236) (0.6800000,-0.2314018) (0.6933333,-0.2357706) (0.7066667,-0.2401299) (0.7200000,-0.2444796) (0.7333333,-0.2488195) (0.7466667,-0.2531493) (0.7600000,-0.2574691) (0.7733333,-0.2617786) (0.7866667,-0.2660777) (0.8000000,-0.2703662) (0.8133333,-0.2746440) (0.8266667,-0.2789110) (0.8400000,-0.2831670) (0.8533333,-0.2874119) (0.8666667,-0.2916455) (0.8800000,-0.2958677) (0.8933333,-0.3000784) (0.9066667,-0.3042774) (0.9200000,-0.3084646) (0.9333333,-0.3126399) (0.9466667,-0.3168032) (0.9600000,-0.3209543) (0.9733333,-0.3250931) (0.9866667,-0.3292195) (1.0000000,-0.3333333) (1.0133333,-0.3374345) (1.0266667,-0.3415229) (1.0400000,-0.3455985) (1.0533333,-0.3496610) (1.0666667,-0.3537105) (1.0800000,-0.3577467) (1.0933333,-0.3617696) (1.1066667,-0.3657790) (1.1200000,-0.3697750) (1.1333333,-0.3737573) (1.1466667,-0.3777259) (1.1600000,-0.3816806) (1.1733333,-0.3856214) (1.1866667,-0.3895482) (1.2000000,-0.3934609) (1.2133333,-0.3973594) (1.2266667,-0.4012436) (1.2400000,-0.4051134) (1.2533333,-0.4089688) (1.2666667,-0.4128096) (1.2800000,-0.4166358) (1.2933333,-0.4204472) (1.3066667,-0.4242439) (1.3200000,-0.4280258) (1.3333333,-0.4317927) (1.3466667,-0.4355446) (1.3600000,-0.4392814) (1.3733333,-0.4430031) (1.3866667,-0.4467096) (1.4000000,-0.4504009) (1.4133333,-0.4540767) (1.4266667,-0.4577372) (1.4400000,-0.4613823) (1.4533333,-0.4650118) (1.4666667,-0.4686258) (1.4800000,-0.4722241) (1.4933333,-0.4758068) (1.5066667,-0.4793738) (1.5200000,-0.4829250) (1.5333333,-0.4864603) (1.5466667,-0.4899798) (1.5600000,-0.4934835) (1.5733333,-0.4969711) (1.5866667,-0.5004428) (1.6000000,-0.5038985) (1.6133333,-0.5073381) (1.6266667,-0.5107617) (1.6400000,-0.5141691) (1.6533333,-0.5175603) (1.6666667,-0.5209354) (1.6800000,-0.5242942) (1.6933333,-0.5276369) (1.7066667,-0.5309632) (1.7200000,-0.5342733) (1.7333333,-0.5375671) (1.7466667,-0.5408445) (1.7600000,-0.5441056) (1.7733333,-0.5473504) (1.7866667,-0.5505788) (1.8000000,-0.5537908) (1.8133333,-0.5569864) (1.8266667,-0.5601656) (1.8400000,-0.5633283) (1.8533333,-0.5664747) (1.8666667,-0.5696046) (1.8800000,-0.5727181) (1.8933333,-0.5758152) (1.9066667,-0.5788958) (1.9200000,-0.5819599) (1.9333333,-0.5850077) (1.9466667,-0.5880390) (1.9600000,-0.5910539) (1.9733333,-0.5940523) (1.9866667,-0.5970344) (2.0000000,-0.6000000) (2.0133333,-0.6029492) (2.0266667,-0.6058821) (2.0400000,-0.6087985) (2.0533333,-0.6116986) (2.0666667,-0.6145824) (2.0800000,-0.6174498) (2.0933333,-0.6203009) (2.1066667,-0.6231357) (2.1200000,-0.6259542) (2.1333333,-0.6287565) (2.1466667,-0.6315425) (2.1600000,-0.6343124) (2.1733333,-0.6370660) (2.1866667,-0.6398035) (2.2000000,-0.6425248) (2.2133333,-0.6452300) (2.2266667,-0.6479191) (2.2400000,-0.6505922) (2.2533333,-0.6532492) (2.2666667,-0.6558903) (2.2800000,-0.6585154) (2.2933333,-0.6611246) (2.3066667,-0.6637178) (2.3200000,-0.6662953) (2.3333333,-0.6688569) (2.3466667,-0.6714027) (2.3600000,-0.6739327) (2.3733333,-0.6764471) (2.3866667,-0.6789458) (2.4000000,-0.6814288) (2.4133333,-0.6838963) (2.4266667,-0.6863482) (2.4400000,-0.6887846) (2.4533333,-0.6912056) (2.4666667,-0.6936111) (2.4800000,-0.6960013) (2.4933333,-0.6983761) (2.5066667,-0.7007357) (2.5200000,-0.7030800) (2.5333333,-0.7054091) (2.5466667,-0.7077232) (2.5600000,-0.7100221) (2.5733333,-0.7123060) (2.5866667,-0.7145749) (2.6000000,-0.7168288) (2.6133333,-0.7190679) (2.6266667,-0.7212922) (2.6400000,-0.7235017) (2.6533333,-0.7256964) (2.6666667,-0.7278765) (2.6800000,-0.7300420) (2.6933333,-0.7321929) (2.7066667,-0.7343293) (2.7200000,-0.7364512) (2.7333333,-0.7385588) (2.7466667,-0.7406520) (2.7600000,-0.7427309) (2.7733333,-0.7447957) (2.7866667,-0.7468462) (2.8000000,-0.7488827) (2.8133333,-0.7509051) (2.8266667,-0.7529135) (2.8400000,-0.7549080) (2.8533333,-0.7568886) (2.8666667,-0.7588554) (2.8800000,-0.7608085) (2.8933333,-0.7627479) (2.9066667,-0.7646737) (2.9200000,-0.7665859) (2.9333333,-0.7684846) (2.9466667,-0.7703698) (2.9600000,-0.7722417) (2.9733333,-0.7741003) (2.9866667,-0.7759456) (3.0000000,-0.7777778) (3.0133333,-0.7795968) (3.0266667,-0.7814028) (3.0400000,-0.7831957) (3.0533333,-0.7849758) (3.0666667,-0.7867429) (3.0800000,-0.7884973) (3.0933333,-0.7902389) (3.1066667,-0.7919679) (3.1200000,-0.7936842) (3.1333333,-0.7953880) (3.1466667,-0.7970794) (3.1600000,-0.7987583) (3.1733333,-0.8004248) (3.1866667,-0.8020791) (3.2000000,-0.8037212) (3.2133333,-0.8053511) (3.2266667,-0.8069689) (3.2400000,-0.8085747) (3.2533333,-0.8101686) (3.2666667,-0.8117505) (3.2800000,-0.8133207) (3.2933333,-0.8148791) (3.3066667,-0.8164257) (3.3200000,-0.8179608) (3.3333333,-0.8194843) (3.3466667,-0.8209963) (3.3600000,-0.8224969) (3.3733333,-0.8239861) (3.3866667,-0.8254640) (3.4000000,-0.8269307) (3.4133333,-0.8283862) (3.4266667,-0.8298306) (3.4400000,-0.8312640) (3.4533333,-0.8326864) (3.4666667,-0.8340979) (3.4800000,-0.8354986) (3.4933333,-0.8368885) (3.5066667,-0.8382677) (3.5200000,-0.8396362) (3.5333333,-0.8409942) (3.5466667,-0.8423416) (3.5600000,-0.8436786) (3.5733333,-0.8450052) (3.5866667,-0.8463215) (3.6000000,-0.8476276) (3.6133333,-0.8489234) (3.6266667,-0.8502091) (3.6400000,-0.8514848) (3.6533333,-0.8527504) (3.6666667,-0.8540062) (3.6800000,-0.8552520) (3.6933333,-0.8564881) (3.7066667,-0.8577144) (3.7200000,-0.8589310) (3.7333333,-0.8601380) (3.7466667,-0.8613354) (3.7600000,-0.8625233) (3.7733333,-0.8637019) (3.7866667,-0.8648710) (3.8000000,-0.8660308) (3.8133333,-0.8671814) (3.8266667,-0.8683228) (3.8400000,-0.8694551) (3.8533333,-0.8705783) (3.8666667,-0.8716926) (3.8800000,-0.8727978) (3.8933333,-0.8738942) (3.9066667,-0.8749818) (3.9200000,-0.8760606) (3.9333333,-0.8771307) (3.9466667,-0.8781922) (3.9600000,-0.8792451) (3.9733333,-0.8802895) (3.9866667,-0.8813254) (4.0000000,-0.8823529)};
\draw[violet!75!black,densely dotted] (-1,-0.3333333)--(1,0.3333333);
\fill[violet!75!black] (-1,-0.3333333) circle(2.3pt);
\fill[violet!75!black] (1,0.3333333) circle(2.3pt);
\end{scope}
\end{scope}
\node[anchor=west,align=left,text width=4.35cm] at (8.9,4.4) {蓝实线：$a=2$\\[8pt]橙虚线：$a=\frac12$\\[14pt]两图都经过原点。\\[8pt]靠近 $y=\pm1$，\\但取不到这两个值。};
\node[text width=12.8cm,align=center] at (7.1,.65) {$D=\mathbb R$；分母始终为正，$f(0)=0$。};
\end{tikzpicture}\end{center}


\textbf{\textcolor{CStatement}{它解决什么问题？}}

遇到 $a^x$、分式、对数、根号混在一起的函数，题目问“奇还是偶”，如果每次都从头计算 $f(-x)$，容易慢，也容易忘记定义域。\textbf{其实许多式子都来自同一件事：把 $x$ 换成 $-x$，有些量变成倒数，有些量保持不变。}

图中 $a=2$ 与 $a=\frac12$ 的曲线都关于原点对称。另一个同类函数是
\[
 f(x)=\frac{2^x-1}{2^x+1},\qquad
 g(x)=\log_2\!\bigl(x+\sqrt{x^2+1}\bigr).
\]
看起来完全不同，却都是奇函数。我们要学会\textbf{认结构、验条件、再速判}。
\medskip
\end{statementbox}
\clearpage
\begin{statementbox}
\textbf{\textcolor{CStatement}{先记共同条件}}

令 $a>0$ 且 $a\ne1$。先检查：\textbf{若 $x$ 能代入，$-x$ 也必须能代入。}然后比较
\[
 \boxed{f(-x)=-f(x)\ (\text{奇})},\qquad
 \boxed{f(-x)=f(x)\ (\text{偶})}.
\]

\textbf{7 个常见奇函数（序号与后文一致）：}
\[
\begin{array}{rl}
\text{奇1}&a^x-a^{-x},\\[3pt]
\text{奇2}&\dfrac{a^x-1}{a^x+1},\\[6pt]
\text{奇3}&\dfrac{a^x+1}{a^x-1}\quad(x\ne0),\\[6pt]
\text{奇4}&\ln\dfrac{kx-b}{kx+b},\\[6pt]
\text{奇5}&\ln\dfrac{kx+b}{kx-b},\\[6pt]
\text{奇6}&\log_a\!\bigl(x+\sqrt{x^2+1}\bigr),\\[3pt]
\text{奇7}&\log_a\!\bigl(\sqrt{x^2+1}-x\bigr).
\end{array}
\]
奇4、奇5 此处取 $k\ne0,b\ne0$，两式均要求 $|x|>|b/k|$；奇6、奇7 在所有实数上有定义。
\end{statementbox}
\clearpage
\begin{statementbox}
\textbf{\textcolor{CStatement}{5 个常见偶函数}}
\[
\begin{array}{rl}
\text{偶1}&a^x+a^{-x},\\[3pt]
\text{偶2}&a^{|x|},\\[3pt]
\text{偶3}&a^{-x^2},\\[3pt]
\text{偶4}&\log_a|x|\quad(x\ne0),\\[3pt]
\text{偶5}&x\left(\dfrac12+\dfrac1{2^x-1}\right)\quad(x\ne0).
\end{array}
\]

\textbf{最快记忆：}
\begin{itemize}[leftmargin=*,itemsep=.25em]
 \item 指数成对：$a^x\pm a^{-x}$，\textbf{加偶、减奇}。
 \item 指数分式：$a^x$ 换成 $a^{-x}=1/a^x$，\textbf{奇2、奇3 都变号}。
 \item 对数比值：换成 $-x$ 时分子分母互换，\textbf{对数变相反数}。
 \item 根式对数：$\sqrt{x^2+1}+x$ 与 $\sqrt{x^2+1}-x$ \textbf{互为倒数}。
 \item 出现 $|x|$ 或 $x^2$，先看是否\textbf{替换后不变}；最后一个偶函数是\textbf{奇函数乘奇函数}。
\end{itemize}

\textbf{边界提醒：}奇函数\textbf{只有在 $0$ 属于定义域时}，才必须满足 $f(0)=0$。奇3、奇4、奇5、偶4、偶5 均不能直接把 $0$ 代入。
\end{statementbox}
